% ===========================================================================
%  00 — 摘要 + Zusammenfassung
%
%  证据：paper_context/claims_registry.md（仅已验证的声明）
%  术语：paper_context/terminology.md
% ===========================================================================

\chapter*{摘要}

许多对安全性和技能要求严苛的领域，包括航空、工业维修、医学以及应急操作，
都要求学习者练习有序规程，其中下一步正确操作取决于当前系统状态。
随着此类训练向高保真沉浸式仿真迁移，教学系统必须能够提供与上下文相关的
帮助，而不能仅依赖于语言模型的流畅性。本论文通过DCS World中的F/A-18C
冷启动规程来研究该问题。论文设计了一个基于证据的教学运行时系统，该系统
结合了规程包状态跟踪、检索到的规程上下文，以及用于从座舱显示屏中提取
结构化视觉事实的微调视觉语言模型（VLM）。两项技术发现支撑了视觉证据层。
第一，8事实本体暴露了将感知与规程完成混合在一起的复合目标所带来的风险。
第二，与运行时对齐的13事实本体，配合小数据LoRA适配，在独立保留集上的
主要Qwen3.5基础模型与LoRA的对比中，将关键误报降低了64--89\%。Qwen3.5-9B
LoRA适配器在928条双语数据行上训练，在两个保留集上达到0.99的事实准确率，
且经验证与训练集零重叠。补充的Qwen3.6和Gemma-4结果表明该方案可应用于
所评估的主干模型，但性能仍与具体主干模型相关。一项包含九名参与者的形成性
VR试点研究提供了描述性证据：所有四名有教学辅助条件的参与者完成了全部
33步规程，而五名无教学辅助的参与者中无人完成；有教学辅助条件下的遗漏
错误被辅助耦合型错误大量取代。更大样本的受控评估仍有待未来工作。

\newpage
\chapter*{Zusammenfassung}

\begin{otherlanguage}{ngerman}
Prozedurales Lernen in hochrealistischer Simulation erfordert das Ausf\"uhren
geordneter, zustandsabh\"angiger Handlungsabfolgen bei gleichzeitiger
\"Uberpr\"ufung von Systemzust\"anden. Herk\"ommliche Trainingshilfen
unterbrechen jedoch entweder die Immersion oder haben keinen Zugriff auf den
Simulatorzustand. Diese Arbeit begegnet dieser L\"ucke durch eine
telemetriegest\"utzte Tutoring-Laufzeitumgebung, die paketbasierte
Zustandsverfolgung, Retrieval-Augmented-LLM-Hilfegenerierung und ein
feinabgestimmtes Vision-Language-Model (VLM) zur strukturierten Extraktion
visueller Fakten aus Cockpit-Anzeigen kombiniert. Zwei technische Ergebnisse
st\"utzen die visuelle Evidenzschicht. Erstens zeigte die 8-Fact-Ontologie das
Risiko zusammengesetzter Ziele, die Wahrnehmung und prozedurale Fertigstellung
vermischen. Zweitens reduzierte die runtime-ausgerichtete 13-Fact-Ontologie in
Verbindung mit kleiner LoRA-Anpassung kritische False-Positives um 64--89\%
im prim\"aren Qwen3.5-Base-gegen-LoRA-Vergleich auf unabh\"angigen Holdout-Sets.
Der Qwen3.5-9B LoRA-Adapter, trainiert mit 928 bilingualen Datenzeilen,
erreicht eine Fact-Genauigkeit von 0,99 auf zwei Holdout-Sets mit verifizierter
Null-\"Uberlappung zum Training. Erg\"anzende Qwen3.6- und Gemma-4-Ergebnisse
zeigen, dass die Rezeptur auf die untersuchten Backbones anwendbar ist; die
Leistung bleibt jedoch backbone-spezifisch. Eine kleine formative
VR-Pilotstudie mit neun Teilnehmern liefert deskriptive
Evidenz: Alle vier Teilnehmer mit Tutor absolvierten die vollst\"andige
33-Schritt-Prozedur, w\"ahrend keiner der f\"unf Teilnehmer ohne Tutor dies
schaffte; Auslassungsfehler wurden in der Tutor-Bedingung weitgehend durch
assistenzgekoppelte Fehler ersetzt. Eine kontrollierte Evaluation mit
gr\"o\ss{}eren Stichproben bleibt zuk\"unftige Arbeit.
\end{otherlanguage}
